% Part V: Worked Examples
% This file is \include'd from main.tex

\chapter{Program Verification}
\label{ch:ex-verification}

We begin Part~V with the domain closest to JuGeo's origins:
program verification.  We walk through the full 4-tuple
$(\STop,\PStr,\CSpec,\GEval)$ for two concrete programs---a
lock-free queue (where obstructions arise) and a sorting function
(where they vanish).

%% ─────────────────────────────────────────────────────────────
\section{Example 1: A Lock-Free Queue}
\label{sec:lfq}

Consider a Michael--Scott lock-free queue with two operations
\texttt{enqueue} and \texttt{dequeue}, each using a
compare-and-swap (CAS) loop.  The correctness goal is
\emph{linearizability}: every concurrent execution is equivalent to
some sequential one.

%% ── 22.1.1  STop ─────────────────────────────────────────────
\subsection{The Semantic Topos $\STop$}

\begin{construction}[Control-flow site]
\label{con:cfg-site}
Let $\mathcal{C}_Q$ be the category whose
\begin{itemize}[nosep]
  \item \textbf{objects} are program points
        $p_0,p_1,\ldots,p_n$ in the control-flow graph (CFG),
  \item \textbf{morphisms} $p_i\to p_j$ are feasible execution paths
        (sequences of basic-block transitions).
\end{itemize}
The Grothendieck topology~$J$ declares a covering family
of a function node $f$ to be the set of its basic blocks:
\[
  \{B_1\hookrightarrow f,\;B_2\hookrightarrow f,\;\ldots\}
  \;\in\; J(f).
\]
\end{construction}

\begin{definition}[Type presheaf]
\label{def:type-presheaf}
The presheaf $\mathcal{F}_{\mathrm{ty}}:\mathcal{C}_Q^{\op}\to\Set$
assigns to each program point~$p$ the set of variable-type bindings
in scope:
\[
  \mathcal{F}_{\mathrm{ty}}(p)
  \;=\;
  \bigl\{(x:\texttt{Node*}),\;(tail:\texttt{AtomicPtr<Node>}),\;\ldots\bigr\}.
\]
Restriction along a morphism $p\to q$ is the usual scope narrowing.
\end{definition}

The \textbf{subobject classifier} $\Omega$ in the topos
$\Sh(\mathcal{C}_Q,J)$ sends each program point~$p$ to the set of
sieves on~$p$.  Concretely, a sieve on~$p$ is a downward-closed set
of paths arriving at~$p$.  The interpretation: a sieve encodes the
set of \emph{reachable} histories at that point.  An assertion
$\varphi$ at~$p$ is therefore a map
$\llbracket\varphi\rrbracket:1\to\Omega$ selecting the sieve of
paths on which $\varphi$ holds.

\begin{example}[Assertions as sieves]
At program point $p_{\mathrm{cas}}$ (the CAS instruction in
\texttt{enqueue}), the assertion ``\texttt{tail} still points to the
last node'' selects the sieve
$S = \{$paths that have not been preempted by another
\texttt{enqueue}$\}$.  This sieve is a \emph{proper} sub-sieve of
the maximal sieve, reflecting the fact that the assertion can be
invalidated by concurrent interference.
\end{example}

%% ── 22.1.2  PStr ─────────────────────────────────────────────
\subsection{The Proof Structure $\PStr$}

\begin{construction}[Specification type]
\label{con:spec-type}
The specification is the type
\[
  \mathrm{Lin}(Q)
  \;:\equiv\;
  \Pi\,(\sigma:\mathrm{ConcExec}(Q)).\;
  \Sigma\,(\tau:\mathrm{SeqExec}(Q)).\;
  \GId_{\gone}\!\bigl(\mathrm{Obs},\;
    \mathrm{obs}(\sigma),\;\mathrm{obs}(\tau)\bigr),
\]
reading: ``for every concurrent execution~$\sigma$, there exists a
sequential execution~$\tau$ with identical observable behaviour.''
\end{construction}

The \textbf{evidence} inhabiting this type is a \emph{simulation
relation} $R\subseteq\mathrm{State}\times\mathrm{AbsState}$ together
with proofs that $R$ is preserved by every atomic step.

\begin{construction}[Hoare-logic proof tree]
\label{con:hoare-tree}
The derivation tree has the shape of a Hoare-logic proof.
For \texttt{enqueue(x)}:
\[
  \dfrac{
    \dfrac{
      \{\texttt{tail}=t\}
      \;\texttt{new\_node := alloc(x)}\;
      \{\texttt{new\_node}\neq\texttt{null}\}
    }{\text{(Alloc)}}
    \quad
    \dfrac{
      \{\texttt{tail}=t \wedge \texttt{next}=\texttt{null}\}
      \;\texttt{CAS(tail.next, null, new\_node)}\;
      \{R(t,t')\}
    }{\text{(CAS-success)}}
  }{
    \{\mathrm{pre}\}\;\texttt{enqueue(x)}\;\{\mathrm{post}\}
  }
\]
Each horizontal line is a typing rule application
$\Gamma\vdash_{\gone} e:A$; the full tree is the term witnessing
$\mathrm{Lin}(Q)$.
\end{construction}

\textbf{$\beta$-reduction.}\quad
In $\PStr$, $\beta$-reduction corresponds to
\emph{proof simplification}.  For instance, composing the CAS-success
lemma with the swing-tail lemma yields a single invariant-preservation
step:
\[
  (\lambda h.\;\mathrm{swing}(h))\;(\mathrm{cas\text{-}ok}(\sigma))
  \;\;\longrightarrow_\beta\;\;
  \mathrm{swing}(\mathrm{cas\text{-}ok}(\sigma)),
\]
reducing two proof steps to one.

%% ── 22.1.3  CSpec ────────────────────────────────────────────
\subsection{The Cohomological Spectrum $\CSpec$}

\begin{construction}[Thread decomposition as a cover]
\label{con:thread-cover}
Partition the CFG into thread-local sub-graphs
$U_{\mathrm{enq}},U_{\mathrm{deq}}$ (one per operation).  This is a
cover in~$J$.  The \v{C}ech complex is:
\[
  C^0 \;=\;
  \mathrm{Proof}(U_{\mathrm{enq}})
  \;\times\;
  \mathrm{Proof}(U_{\mathrm{deq}}),
  \qquad
  C^1 \;=\;
  \mathrm{Proof}(U_{\mathrm{enq}}\cap U_{\mathrm{deq}}).
\]
Here $C^0$ contains each thread's \emph{local} correctness proof,
and $C^1$ contains proofs about shared state accessed by both threads.
\end{construction}

The coboundary $\delta^0:C^0\to C^1$ computes:
\[
  \delta^0(s_{\mathrm{enq}},s_{\mathrm{deq}})
  \;=\;
  s_{\mathrm{enq}}\big|_{U_{\mathrm{enq}}\cap U_{\mathrm{deq}}}
  \;-\;
  s_{\mathrm{deq}}\big|_{U_{\mathrm{enq}}\cap U_{\mathrm{deq}}}.
\]

\begin{proposition}[Obstruction = data race]
\label{prop:race-obstruction}
$H^1(\mathcal{C}_Q,\mathcal{F}_{\mathrm{inv}})\neq 0$ if and only
if the local invariants of \texttt{enqueue} and \texttt{dequeue}
are inconsistent on shared state---i.e., there exists a data race
or an atomicity violation.
\end{proposition}

\begin{example}[A concrete race]
Suppose a buggy implementation omits the ``swing tail'' step after
a successful CAS.  Then \texttt{enqueue} believes
$\texttt{tail}=\texttt{new\_node}$ on exit, while a concurrent
\texttt{dequeue} still sees the old tail.  The difference
\[
  \delta^0(\texttt{tail}=\texttt{new\_node},\;
           \texttt{tail}=\texttt{old\_tail})
  \;\neq\;0
\]
is a non-trivial 1-cocycle: the obstruction.
Mayer--Vietoris gives the exact sequence
\[
  \cdots\to
  H^0(U_{\mathrm{enq}})\oplus H^0(U_{\mathrm{deq}})
  \xrightarrow{\;\delta^0\;}
  H^0(U_{\mathrm{enq}}\cap U_{\mathrm{deq}})
  \to
  H^1(\mathcal{C}_Q,\mathcal{F}_{\mathrm{inv}})
  \to\cdots
\]
The non-vanishing $H^1$ localises the bug to the shared-state
interface.
\end{example}

\textbf{Repair.}\quad
Adding the swing-tail step makes $\delta^0=0$, whence
$H^1=0$ and the proofs glue into a global linearizability argument.

%% ── 22.1.4  GEval ────────────────────────────────────────────
\subsection{The Graded Evaluation $\GEval$}

\begin{construction}[Verification grade]
\label{con:verif-grade}
The overall grade is a product in the grade semiring:
\[
  G
  \;=\;
  g_{\mathrm{test}} \;\otG\; g_{\mathrm{type}} \;\otG\; g_{\mathrm{perf}},
\]
where each component lies in $\Grade=(-\infty,0]\cup\{-\infty\}$
(with $\gone=0$ perfect, $\gzero=-\infty$ vacuous).
\end{construction}

\begin{itemize}[nosep]
  \item \textbf{Test-coverage grade} $g_{\mathrm{test}}$.
    The trail is: unit tests ($g=-2.1$, covering basic enqueue/dequeue)
    $\to$ integration tests ($g=-1.3$, two-thread interleaving)
    $\to$ stress tests ($g=-0.4$, 64 threads, 10\textsuperscript{6}
    operations).  Convergence of random fuzzing pushes $g_{\mathrm{test}}$
    toward~$\gone$.

  \item \textbf{Type-safety grade} $g_{\mathrm{type}}$.
    With a verified simulation relation checked by a proof assistant,
    $g_{\mathrm{type}}=\gone$ (the proof is machine-checked).

  \item \textbf{Performance grade} $g_{\mathrm{perf}}$.
    Benchmarks show lock-free throughput exceeds a mutex-based queue
    by $3\times$ at 32~threads.  We assign $g_{\mathrm{perf}}=-0.2$
    (minor contention penalty at very high thread counts).
\end{itemize}

The overall grade is
$G = (-0.4)\otG 0 \otG (-0.2) = -0.4$
(bottleneck: test coverage not yet exhaustive).

\medskip

%% ─────────────────────────────────────────────────────────────
\section{Example 2: Insertion Sort}
\label{sec:isort}

To contrast, consider a sequential insertion sort on integer arrays.

\subsection{The 4-Tuple in Brief}

\textbf{$\STop$.}\quad
The site $\mathcal{C}_S$ has program points
$\{p_0,\ldots,p_5\}$ with linear control flow (one loop).  The
topology is trivial: the single function is its own cover.
The type presheaf assigns $\texttt{int[]}$ at every point.

\textbf{$\PStr$.}\quad
The specification type is
$\mathrm{Sorted}(a') \;\wedge\; \mathrm{Perm}(a,a')$.
The proof is a standard loop-invariant argument:
\[
  \dfrac{
    \Gamma\vdash_{\gone}
      \mathrm{inv}(i):\mathrm{Sorted}(a[0..i])
    \quad
    \Gamma\vdash_{\gone}
      \mathrm{step}(i):\mathrm{inv}(i)\Rightarrow\mathrm{inv}(i+1)
  }{
    \Gamma\vdash_{\gone}
      \mathrm{loop}:\mathrm{Sorted}(a[0..n])
  }
\]
$\beta$-reduction unfolds the loop to $n$ applications of
$\mathrm{step}$.

\textbf{$\CSpec$.}\quad
Since the program is sequential, the cover has a single element
(the whole function).  The \v{C}ech complex degenerates:
$C^0=\mathrm{Proof}(U)$, $C^1=0$.  Therefore $H^1=0$---no
obstructions.  The proof glues trivially.

\textbf{$\GEval$.}\quad
$g_{\mathrm{test}}=-0.3$ (exhaustive tests up to $n=10\,000$),
$g_{\mathrm{type}}=\gone$ (mechanised proof of the loop invariant),
$g_{\mathrm{perf}}=-1.5$ ($O(n^2)$ is slow for large~$n$).
Overall $G=-1.5$ (bottleneck: performance).

\begin{remark}[Comparing the two examples]
The lock-free queue has non-trivial $H^1$ that must be resolved;
insertion sort has $H^1=0$ by triviality of the cover.  Both
achieve high type-safety grades, but differ in their test-coverage
and performance profiles.  The 4-tuple captures these differences
precisely: the geometric (cohomological) and evaluative aspects
are orthogonal dimensions of program quality.
\end{remark}

\chapter{Natural Language Semantics}
\label{ch:ex-nl}

We now turn to natural language, the domain where the CTT extension
of JuGeo (Part~III) is most directly exercised.  Three examples
illustrate the 4-tuple at progressively greater interpretive depth:
a classic donkey sentence, a garden-path sentence, and a metaphor.

%% ─────────────────────────────────────────────────────────────
\section{Example 1: Donkey Anaphora}
\label{sec:donkey}

\begin{quote}
\emph{``Every farmer who owns a donkey beats it.''}
\end{quote}

This sentence is the touchstone of formal semantics: the pronoun
\emph{it} must be bound by the indefinite \emph{a donkey}, yet
the two sit in different clauses.  We show how each of the four
objects handles this.

%% ── STop ─────────────────────────────────────────────────────
\subsection{The Semantic Topos $\STop$}

\begin{construction}[NL site]
\label{con:nl-site}
The category $\mathcal{C}_{\mathrm{NL}}$ has:
\begin{itemize}[nosep]
  \item \textbf{Objects}: discourse segments at every granularity---morphemes,
    words, phrases, clauses, sentences, and multi-sentence discourses.
  \item \textbf{Morphisms}: inclusion of a sub-segment into a
    larger segment (e.g.\ the relative clause ``who owns a donkey''
    into the full sentence).
\end{itemize}
The topology~$J$ declares a covering family of a sentence to
be its immediate-constituent phrases:
\[
  \bigl\{\text{NP}_{\text{subj}}\hookrightarrow S,\;
         \text{VP}\hookrightarrow S\bigr\}
  \;\in\;J(S).
\]
\end{construction}

\begin{definition}[Semantic-type presheaf]
The presheaf $\mathcal{F}_{\sem}:\mathcal{C}_{\mathrm{NL}}^{\op}\to\Set$
assigns to each node its semantic type in the stratal universe tower:
\[
  \mathcal{F}_{\sem}(\text{``farmer''}) = U_{\sem},\quad
  \mathcal{F}_{\sem}(\text{``beats''}) = U_{\sem}\to U_{\sem}\to U_{\sem},\quad
  \mathcal{F}_{\sem}(\text{``it''}) = U_{\disc}.
\]
The pronoun lives at the \emph{discourse} stratum $U_{\disc}$
because its referent is determined by cross-clausal context.
\end{definition}

The internal logic of $\Sh(\mathcal{C}_{\mathrm{NL}},J)$ is
intuitionistic: truth of a proposition at a discourse segment~$d$
means truth in \emph{every extension} of~$d$.  This captures the
open-world nature of discourse: adding more sentences can never
revoke what is established.

%% ── PStr ─────────────────────────────────────────────────────
\subsection{The Proof Structure $\PStr$}

We build the semantic representation step by step.

\begin{construction}[Lexicon entries as types]
\label{con:lexicon}
\begin{align*}
  \text{every}   &:\; \Pi(P:U_{\sem}\to\Type).\,
                       \Pi(Q:U_{\sem}\to\Type).\,
                       \bigl(\Pi(x:U_{\sem}).\,P(x)\Rightarrow Q(x)\bigr), \\
  \text{farmer}  &:\; U_{\sem}\to\Prop, \\
  \text{owns}    &:\; U_{\sem}\to U_{\sem}\to\Prop, \\
  \text{a}       &:\; \Pi(P:U_{\sem}\to\Type).\,\Sigma(x:U_{\sem}).\,P(x), \\
  \text{donkey}  &:\; U_{\sem}\to\Prop, \\
  \text{beats}   &:\; U_{\sem}\to U_{\sem}\to\Prop, \\
  \text{it}      &:\; \text{(unresolved discourse referent)}.
\end{align*}
\end{construction}

\textbf{The anaphora problem.}\quad
Under na\"ive composition, ``a donkey'' introduces an existential
$\Sigma(y:U_{\sem}).\,\mathrm{donkey}(y)$ inside the relative clause,
but the pronoun ``it'' sits in the matrix clause and must refer
to~$y$.  The variable~$y$ is not in scope.

\begin{construction}[Resolution via the discourse monad]
\label{con:dmon}
Define the discourse monad
$\mathrm{DMon}(A) = \mathrm{DCtx}\to(A\times\mathrm{DCtx})$
where $\mathrm{DCtx}$ is a list of discourse referents.
The indefinite ``a donkey'' performs a monadic action that
\emph{pushes} a referent~$y$ onto $\mathrm{DCtx}$:
\[
  \llbracket\text{a donkey}\rrbracket
  \;=\;
  \lambda\,\gamma.\;
  \mathrm{let}\;y=\mathrm{fresh}(\gamma)\;\mathrm{in}\;
  (y,\;\gamma\!::\!y).
\]
The pronoun ``it'' is a \emph{lookup}:
$\llbracket\text{it}\rrbracket = \lambda\gamma.\,
  (\mathrm{top}(\gamma),\;\gamma)$.
\end{construction}

Composing inside $\mathrm{DMon}$, the full derivation is:
\[
  \dfrac{
    \dfrac{
      \Gamma\vdash_{\gone} \text{farmer}:U_{\sem}\to\Prop
      \qquad
      \Gamma\vdash_{\gone}
        [\text{who owns a donkey}] :
        \Pi(x).\,\mathrm{DMon}(\mathrm{donkey}(y)\wedge\mathrm{owns}(x,y))
    }{
      \Gamma\vdash_{\gone}
        [\text{farmer who owns a donkey}] :
        \Pi(x).\,\mathrm{DMon}(\mathrm{farmer}(x)\wedge
                                \mathrm{donkey}(y)\wedge
                                \mathrm{owns}(x,y))
    }
  }{
    \Gamma\vdash_{\gone}
      S : \Pi(x).\,
          \bigl(\mathrm{farmer}(x)\wedge\mathrm{donkey}(y)
                \wedge\mathrm{owns}(x,y)\bigr)
          \Rightarrow\mathrm{beats}(x,y)
  }
\]
The monadic plumbing threads $y$ from the relative clause to the
matrix clause, binding ``it'' to the donkey.

%% ── CSpec ────────────────────────────────────────────────────
\subsection{The Cohomological Spectrum $\CSpec$}

\begin{construction}[Clause cover]
\label{con:clause-cover}
Take the cover $\{U_{\mathrm{rel}},U_{\mathrm{mat}}\}$ where
$U_{\mathrm{rel}}$ = the relative clause and
$U_{\mathrm{mat}}$ = the matrix clause.
\end{construction}

\textbf{Before resolution.}\quad
The presheaf of discourse referents restricted to $U_{\mathrm{rel}}$
contains~$y$, but restricted to $U_{\mathrm{mat}}$ does not.
The coboundary is:
\[
  \delta^0(s_{\mathrm{rel}},s_{\mathrm{mat}})
  \;=\;
  s_{\mathrm{rel}}\big|_{U_{\mathrm{rel}}\cap U_{\mathrm{mat}}}
  - s_{\mathrm{mat}}\big|_{U_{\mathrm{rel}}\cap U_{\mathrm{mat}}}
  \;\neq\; 0,
\]
since $y\in s_{\mathrm{rel}}$ but $y\notin s_{\mathrm{mat}}$.
Therefore $H^1\neq 0$: the anaphora is an obstruction to local
gluing.

\textbf{After resolution.}\quad
Switching to the discourse-monad presheaf (which threads $y$
through $\mathrm{DCtx}$) makes $y$ available at
$U_{\mathrm{mat}}$.  Now $\delta^0=0$ and $H^1=0$: the
interpretation glues globally.

\begin{remark}
This is the cohomological content of Dynamic Predicate Logic
and DRT: the discourse monad is exactly the change-of-coefficients
that kills $H^1$.
\end{remark}

%% ── GEval ────────────────────────────────────────────────────
\subsection{The Graded Evaluation $\GEval$}

The sentence receives independent grades at each linguistic stratum:
\begin{align*}
  g_{\phon}  &= -0.05  &&\text{(phonologically well-formed, minor coda cluster)},\\
  g_{\synt}  &= \gone   &&\text{(grammatical)},\\
  g_{\sem}   &= -0.1   &&\text{(coherent once anaphora is resolved)},\\
  g_{\prag}  &= -0.3   &&\text{(pragmatically odd: ``beats'' is harsh)}.
\end{align*}
The overall grade is
$T = g_{\phon}\otG g_{\synt}\otG g_{\sem}\otG g_{\prag}
   = -0.3$
(bottleneck: pragmatic felicity).
In the normalised scale where $\gone=0$ is perfect, a grade of
$-0.3$ corresponds to ``well-formed but stylistically marked.''

\medskip

%% ─────────────────────────────────────────────────────────────
\section{Example 2: Garden-Path Sentence}
\label{sec:garden-path}

\begin{quote}
\emph{``The horse raced past the barn fell.''}
\end{quote}

This sentence is grammatical (a reduced relative clause:
``The horse [that was] raced past the barn fell''),
but readers initially parse ``raced'' as the main verb and
reach ``fell'' with no syntactic slot for it.

\textbf{$\STop$.}\quad
The site is the same NL category $\mathcal{C}_{\mathrm{NL}}$.
The initial parse assigns the cover
$\{$NP:``the horse'', VP:``raced past the barn''$\}$
to the substring before ``fell.''

\textbf{$\PStr$ (first attempt).}\quad
The first derivation tree types ``raced'' as a main verb:
$\Gamma\vdash_{\gone}\text{raced}:U_{\sem}\to U_{\sem}$.
When ``fell'' arrives, the context $\Gamma$ has no open
argument position.  The typing derivation fails:
$\Gamma\nvdash\text{fell}:A$ for any~$A$.

\textbf{$\CSpec$ (first attempt).}\quad
The cover $\{U_{\text{pre-fell}}, U_{\text{fell}}\}$ has
$H^1\neq 0$: the local parse of ``raced past the barn'' does not
extend to include ``fell.''  The obstruction is localised at the
word boundary before ``fell.''

\textbf{Reanalysis.}\quad
The parser \emph{recovers} the cover: ``raced past the barn'' is
re-typed as a reduced relative clause (participial modifier), freeing
a VP slot for ``fell.''
Under the new cover,
$\Gamma\vdash_{\gone}[\text{the horse raced past the barn}]:
  \mathrm{NP}$
and
$\Gamma\vdash_{\gone}\text{fell}:\mathrm{NP}\to S$,
giving $H^1=0$.

\textbf{$\GEval$.}\quad
After reanalysis:
$g_{\synt}=-0.6$ (syntactically legal but highly marked),
$g_{\sem}=\gone$ (semantically clear),
$g_{\prag}=-0.8$ (pragmatically confusing).
Overall $T=-0.8$: a difficult sentence.

\medskip

%% ─────────────────────────────────────────────────────────────
\section{Example 3: Metaphor}
\label{sec:metaphor}

\begin{quote}
\emph{``Juliet is the sun.''}  --- Shakespeare, \emph{Romeo and Juliet}
\end{quote}

\textbf{$\STop$.}\quad
The site is $\mathcal{C}_{\mathrm{NL}}$; the relevant presheaf
is $\mathcal{F}_{\sem}$ assigning semantic types.  At the literal
level, $\mathcal{F}_{\sem}(\text{Juliet})=\texttt{Person}$ and
$\mathcal{F}_{\sem}(\text{sun})=\texttt{Star}$.

\textbf{$\PStr$.}\quad
A literal identity $\GId_{\gone}(\texttt{Entity},\text{Juliet},\text{sun})$
is uninhabited: a person is not a star.  But a \emph{graded} identity
\[
  \GId_g(\texttt{Entity},\text{Juliet},\text{sun})
  \quad\text{with}\quad g\approx -0.35
\]
is inhabited when we restrict to the sub-type of entities possessing
\emph{radiance, warmth, centrality}.  The grade $g$ measures the
proportion of shared structure.

The typing derivation is:
\[
  \dfrac{
    \Gamma\vdash_{\gone}\text{Juliet}:\texttt{Person}
    \qquad
    \Gamma\vdash_{\gone}\text{sun}:\texttt{Star}
    \qquad
    \Gamma\vdash_{-0.35}
      \mathrm{radiance}:\texttt{Person}\cap\texttt{Star}
  }{
    \Gamma\vdash_{-0.35}
      \text{``Juliet is the sun''}
      :\GId_{-0.35}(\texttt{Entity},\text{Juliet},\text{sun})
  }
\]

\textbf{$\CSpec$.}\quad
$H^1=0$ at the metaphorical grade: the identification glues
across the source and target domains because we have restricted
to shared features.  At the literal grade ($g=\gone$), $H^1\neq 0$:
the literal identification fails.

\textbf{$\GEval$.}\quad
$g_{\phon}=\gone$ (simple monosyllables),
$g_{\synt}=\gone$ (subject--copula--predicate),
$g_{\sem}=-0.35$ (metaphorical, not literal),
$g_{\poet}=-0.05$ (vivid imagery, canonical).
The \emph{poetic} grade is near-perfect; the semantic grade reflects
the deliberate non-literalness.  Overall
$T=-0.35$ in the semantic dimension, but the evaluative judgment
is that this is \emph{excellent} poetry: the aesthetic grade overrides
the semantic gap.

\begin{remark}[Graded identity as the algebra of figurative language]
The metaphor example illustrates the central thesis of CTT:
figurative language is not a ``deviation'' from literal meaning
but a principled use of the graded identity type $\GId_g$ at a
grade $\gzero < g < \gone$.  Simile, metonymy, synecdoche, and
irony each correspond to different structural restrictions on
the shared sub-type that witnesses the graded identification.
\end{remark}

\chapter{Poetry Generation and Analysis}
\label{ch:ex-poetry}

Poetry is the domain where all four objects are most densely
populated.  We present three examples: generating a Shakespearean
sonnet, analysing a Keats sonnet, and diagnosing a flawed poem.

%% ─────────────────────────────────────────────────────────────
\section{Example 1: Generating a Shakespearean Sonnet}
\label{sec:sonnet-gen}

\textbf{Prompt.}\; Generate a Shakespearean sonnet on
\emph{the passage of time}.

%% ── STop ─────────────────────────────────────────────────────
\subsection{The Semantic Topos $\STop$}

\begin{construction}[Poetic site]
\label{con:poet-site}
The category $\mathcal{C}_{\mathrm{poet}}$ refines
$\mathcal{C}_{\mathrm{NL}}$ with poetic structure:
\begin{itemize}[nosep]
  \item \textbf{Objects}: syllables, metrical feet, lines,
    quatrains (Q1--Q3), couplet (C), and the full sonnet~$S$.
  \item \textbf{Morphisms}: inclusions
    $\text{syll}\hookrightarrow\text{foot}\hookrightarrow
     \text{line}\hookrightarrow\text{quatrain}\hookrightarrow S$.
\end{itemize}
The topology declares the 14-line decomposition as a cover of~$S$:
\[
  \{l_1,l_2,\ldots,l_{14}\hookrightarrow S\}\in J(S),
\]
and at a coarser level,
$\{Q_1,Q_2,Q_3,C\hookrightarrow S\}\in J(S)$.
\end{construction}

The presheaf $\mathcal{F}_{\mathrm{ling}}$ assigns to each level
the available linguistic material: at the syllable level, phonemes
from the CMU dictionary; at the line level, syntactic frames and
semantic content; at the quatrain level, rhetorical functions
(exposition, complication, turn, resolution).

%% ── PStr ─────────────────────────────────────────────────────
\subsection{The Proof Structure $\PStr$}

\begin{construction}[Sonnet type]
\label{con:sonnet-type}
The specification type is:
\[
  \mathrm{Sonnet}(\theta, \rho, \mu)
  \;:\equiv\;
  \Sigma(l_{1..14}:\mathrm{Line}^{14}).\;
  \mathrm{Theme}(l_{1..14},\theta)
  \;\wedge\;
  \mathrm{Rhyme}(l_{1..14},\rho)
  \;\wedge\;
  \mathrm{Meter}(l_{1..14},\mu),
\]
where $\theta=\text{time}$,
$\rho=\text{ABAB\,CDCD\,EFEF\,GG}$,
$\mu=\text{iambic pentameter}$.
\end{construction}

\textbf{Step 1: Choose a conceit.}\quad
The $\FMLM$ proposes three conceits---time as \emph{river},
\emph{season}, \emph{hourglass}---and scores each against the
theme type.  ``River'' wins ($g=-0.05$) for richness of imagery.

\textbf{Step 2: Generate line by line.}\quad
Each line~$l_i$ is typed in context:
\[
  \Gamma,\;l_1,\ldots,l_{i-1}
  \;\vdash_{g_i}\;
  l_i : \mathrm{Line}(\theta_i,\rho_i,\mu),
\]
where $\theta_i$ is the local sub-theme and $\rho_i$ constrains the
end-rhyme.  The $\FMLM$ fills each line via masked-slot iteration:
\begin{enumerate}[nosep]
  \item Start with a template:
    \texttt{[\_\_\_] of time [\_\_\_] like a [\_\_\_]}.
  \item Generate candidates for each slot from WordNet/FrameNet.
  \item Score: $g_{ij}=\mathrm{grade}(\Gamma\vdash l_i[w_j/M_k]:T_k(w_j))$.
  \item Select $\arg\max_j\,g_{ij}$.
  \item Iterate until grades stabilise.
\end{enumerate}

\textbf{Step 3: Meter check.}\quad
Each line must satisfy $\mathrm{Meter}(l_i,\mu)$: exactly 10
syllables with stress pattern $\cup/\cup/\cup/\cup/\cup/$ (iambic
pentameter).  The $\HPLF$ voice bundle at the $\phon$ stratum
verifies this:
\[
  G_i^{\phon}(l_i)
  \;=\;
  \gone \;\text{if stress pattern matches},
  \qquad
  g < \gone\;\text{otherwise (substitutions allowed at a cost)}.
\]
A feminine ending (11th unstressed syllable) costs $g=-0.1$.

\textbf{Step 4: Rhyme check.}\quad
End-words of paired lines must satisfy
$\GId_{g_r}(U_{\phon},\mathrm{coda}(w_i),\mathrm{coda}(w_j))$
with $g_r$ close to $\gone$.  Perfect rhyme: $g_r=\gone$.
Near-rhyme (e.g.\ ``love''/``prove''): $g_r\approx -0.15$.

\medskip
Suppose the $\FMLM$ produces the opening quatrain:
\begin{verse}
  The river of our years flows ever on,\\
  Through meadows green where youthful pleasures play,\\
  Until the golden afternoon is gone\\
  And twilight steals the colours from the day.
\end{verse}
Each line has been typed at grade $g_i\in[-0.15,\,0]$; the quatrain
as a whole satisfies $\mathrm{Rhyme}(Q_1,\text{ABAB})$ and
$\mathrm{Meter}(Q_1,\mu)$.

%% ── CSpec ────────────────────────────────────────────────────
\subsection{The Cohomological Spectrum $\CSpec$}

With the quatrain cover $\{Q_1,Q_2,Q_3,C\}$:

\begin{example}[Obstruction at the Q2--Q3 boundary]
Suppose line~5 (start of Q2) rhymes correctly with line~7 but
introduces the word ``frozen,'' breaking the river metaphor
(rivers don't freeze in this conceit).  The \v{C}ech coboundary
between Q1 and Q2 is:
\[
  \delta^0(s_{Q_1},s_{Q_2})
  \;=\;
  \text{``river (flowing)''}\big|_{Q_1\cap Q_2}
  \;-\;
  \text{``frozen (static)''}\big|_{Q_1\cap Q_2}
  \;\neq\;0.
\]
$H^1\neq 0$: the metaphor does not glue across quatrains.

\textbf{Repair}: replace ``frozen'' with ``narrowing'' (the
river narrows as time shrinks).  Now $\delta^0=0$ and
$H^1=0$.
\end{example}

\begin{example}[Volta obstruction]
The volta at line~9 must shift the argument (from ``time
passes'' to ``but something endures'').  If Q3 merely
continues the lament without a turn, the structural
constraint $\mathrm{Volta}(l_9)$ is unsatisfied---an
obstruction at the $Q_2$--$Q_3$ boundary in the rhetorical
presheaf.  Repair: rewrite line~9 as a ``But'' or ``Yet''
clause.
\end{example}

%% ── GEval ────────────────────────────────────────────────────
\subsection{The Graded Evaluation $\GEval$}

After repair, the final grades (in the $\HPLF$ framework) are:
\begin{align*}
  g_{\mathrm{meter}}   &= -0.08  &&\text{(one feminine ending in line 12)},\\
  g_{\mathrm{rhyme}}   &= -0.05  &&\text{(one near-rhyme: ``on''/``gone'')},\\
  g_{\mathrm{imagery}} &= -0.12  &&\text{(one clich\'e: ``golden afternoon'')},\\
  g_{\mathrm{arc}}     &= -0.10  &&\text{(volta could be sharper)}.
\end{align*}
The overall harmony:
\[
  \mathrm{Harm}^{\HPLF}
  \;=\;
  g_{\mathrm{meter}} \otG g_{\mathrm{rhyme}} \otG
  g_{\mathrm{imagery}} \otG g_{\mathrm{arc}}
  \;=\; -0.12
\]
(bottleneck: imagery).  The sonnet is competent but not
outstanding; the clich\'e drags the grade down.

\medskip

%% ─────────────────────────────────────────────────────────────
\section{Example 2: Analysing a Keats Sonnet}
\label{sec:keats}

\begin{quote}
\emph{``When I have fears that I may cease to be \\
Before my pen has glean'd my teeming brain\ldots''}\\
\hfill --- Keats, \emph{Sonnet VII}
\end{quote}

\textbf{$\STop$.}\quad
The poetic site is the same $\mathcal{C}_{\mathrm{poet}}$ as above,
with the 14-line cover.  The presheaf assigns Keats's actual
linguistic material at each node.

\textbf{$\PStr$.}\quad
The type $\mathrm{Sonnet}(\text{mortality},\text{ABAB}\ldots,\mu)$
is inhabited by the actual poem.  Key typing observations:
\begin{itemize}[nosep]
  \item ``Glean'd my teeming brain'' uses $\GId_{-0.2}$
    (metaphor: mind = field to be harvested).
  \item ``High-pil\`ed books'' carries
    $\mathrm{Ico}(\text{books},\text{grain})$: an iconicity between
    physical stacking and intellectual accumulation.
  \item The volta at line~9 (``And when I feel\ldots'') introduces
    the beloved, shifting from self to other.
\end{itemize}

\textbf{$\CSpec$.}\quad
$H^1=0$ throughout: every metaphor, every rhyme, every metrical
choice glues perfectly.  The \v{C}ech complex is exact---the
hallmark of a masterwork.

\textbf{$\GEval$.}\quad
$g_{\mathrm{meter}}=\gone$ (flawless iambic pentameter),
$g_{\mathrm{rhyme}}=\gone$ (all perfect rhymes),
$g_{\mathrm{imagery}}=-0.02$ (near-perfect; every image is
fresh),
$g_{\mathrm{arc}}=-0.03$ (the volta is slightly abrupt).
$\mathrm{Harm}^{\HPLF}=-0.03$.  This is essentially a
grade-$\gone$ poem.

\medskip

%% ─────────────────────────────────────────────────────────────
\section{Example 3: A Flawed Poem}
\label{sec:bad-poem}

Consider the following attempt:
\begin{verse}
  The clock goes tick and also tock,\\
  It tells the time just like a clock,\\
  The hours pass both fast and slow,\\
  And that is all I need to know.
\end{verse}

\textbf{$\STop$.}\quad
The site has 4 lines; the cover is $\{l_1,l_2,l_3,l_4\}$.

\textbf{$\PStr$.}\quad
The type is $\mathrm{Quatrain}(\text{time},\text{AABB},\mu')$
for some meter~$\mu'$.  The poem inhabits this type at a low
grade:
\begin{itemize}[nosep]
  \item Line~2 is a tautology (``like a clock'') --- the semantic
    content is $\GId_{\gone}(\text{clock},\text{clock})$, which
    is trivial and adds no information.
  \item Line~4 is a clich\'e (``all I need to know'').
\end{itemize}

\textbf{$\CSpec$.}\quad
$H^1\neq 0$ at the $l_1$--$l_2$ boundary: line~2 repeats
the content of line~1 without advancing the argument.  The
coboundary
$\delta^0(s_1,s_2) = \text{clock}-\text{clock}$ would be zero
\emph{semantically}, but in the poetic presheaf (which values
\emph{novelty}), the redundancy registers as a non-trivial
cocycle.  Similarly, $H^1\neq 0$ at $l_3$--$l_4$ (clich\'e is
an obstruction to freshness).

\textbf{$\GEval$.}\quad
$g_{\mathrm{meter}}=-0.3$ (irregular stress),
$g_{\mathrm{rhyme}}=-0.5$ (``clock''/``clock'' is an identical
rhyme, penalised),
$g_{\mathrm{imagery}}=-2.0$ (no fresh imagery),
$g_{\mathrm{arc}}=-1.5$ (no development).
$\mathrm{Harm}^{\HPLF}=-2.0$ (bottleneck: imagery).
The poem is essentially at the $\gzero$ end of the scale.

\begin{remark}[Diagnosis vs.\ generation]
The 4-tuple does not merely \emph{judge} the bad poem; it
\emph{localises} the failures.  $\CSpec$ points to the
$l_1$--$l_2$ boundary and the $l_3$--$l_4$ boundary;
$\GEval$ identifies imagery as the worst dimension.  A
generation system could use this diagnosis to \emph{repair}
the poem: replace line~2 with a non-tautological elaboration,
replace line~4's clich\'e, and re-check $H^1$.
\end{remark}

\chapter{Mathematical Ideation and Proof}
\label{ch:ex-math}

Mathematics is the domain where $\PStr$ is richest and $\GEval$
most contested.  We work through the irrationality of $\sqrt{2}$
in full 4-tuple detail, then sketch mathematical ideation as
proof search.

%% ─────────────────────────────────────────────────────────────
\section{Example 1: $\sqrt{2}$ Is Irrational}
\label{sec:sqrt2}

\subsection{The Semantic Topos $\STop$}

\begin{construction}[Mathematical site]
\label{con:math-site}
The category $\mathcal{C}_{\mathrm{math}}$ has:
\begin{itemize}[nosep]
  \item \textbf{Objects}: mathematical statements---definitions
    ($\mathrm{gcd}$, ``rational'', ``irrational''), lemmas
    (``if $p^2$ is even then $p$ is even''), and theorems
    ($\sqrt{2}\notin\mathbb{Q}$).
  \item \textbf{Morphisms}: logical dependency.  A morphism
    $L\to T$ exists when lemma~$L$ is used in the proof of
    theorem~$T$.
\end{itemize}
The topology declares a cover of a theorem to be a
\emph{complete} set of dependencies:
\[
  \bigl\{\text{Def}(\gcd),\;
         \text{Lem}(\text{even-square}),\;
         \text{Ax}(\mathbb{Z}\text{-domain})
  \bigr\}
  \;\in\;J\bigl(\sqrt{2}\notin\mathbb{Q}\bigr).
\]
\end{construction}

The presheaf $\mathcal{F}_{\mathrm{obj}}$ assigns to each node the
set of mathematical objects in scope: at $\text{Def}(\gcd)$, the
objects $\{p,q,\gcd(p,q)\}$; at the theorem node, the full
collection $\{p,q,k,\gcd,\sqrt{2}\}$.

\subsection{The Proof Structure $\PStr$}

\begin{construction}[Irrationality type]
The specification is:
\[
  \mathrm{Irrational}(\sqrt{2})
  \;:\equiv\;
  \neg\,\Sigma(p:\mathbb{Z}).\,\Sigma(q:\mathbb{Z}_{>0}).\,
  \GId_{\gone}\!\bigl(\mathbb{R},\;\sqrt{2},\;p/q\bigr).
\]
\end{construction}

The proof term inhabiting this type is a \emph{proof by
contradiction}.  We present the full derivation tree.

\begin{construction}[Proof tree]
\label{con:sqrt2-proof}
\[
\dfrac{
  \dfrac{
    \dfrac{
      [h:\sqrt{2}=p/q,\;\gcd(p,q)=1]
    }{
      2q^2=p^2
    }\text{\scriptsize(algebra)}
    \qquad
    \dfrac{
      2q^2=p^2
    }{
      2\mid p^2
    }\text{\scriptsize(div)}
    \qquad
    \dfrac{
      2\mid p^2
    }{
      2\mid p
    }\text{\scriptsize(prime)}
  }{
    \Gamma,h\vdash_{\gone} e_1:2\mid p
  }
  \quad
  \dfrac{
    \dfrac{
      p=2k
    }{
      2q^2=4k^2
    }\text{\scriptsize(subst)}
    \quad
    \dfrac{
      q^2=2k^2
    }{
      2\mid q
    }\text{\scriptsize(same)}
  }{
    \Gamma,h\vdash_{\gone} e_2:2\mid q
  }
}{
  \dfrac{
    \Gamma,h\vdash_{\gone} (e_1,e_2):2\mid p\;\wedge\;2\mid q
    \qquad
    \Gamma\vdash_{\gone} c:\gcd(p,q)=1
  }{
    \Gamma\vdash_{\gone}\;\bot
  }\text{\scriptsize(contra)}
}
\]
The overall term is
$\lambda h.\;\mathrm{absurd}(\mathrm{gcd\text{-}contra}(e_1,e_2,c))
 :\mathrm{Irrational}(\sqrt{2})$.
\end{construction}

\textbf{$\beta$-reduction.}\quad
Inlining $k=p/2$ into the second branch and simplifying
$2q^2=4(p/2)^2\longrightarrow q^2=2(p/2)^2$ is a $\beta$-step
that shortens the derivation.

\subsection{The Cohomological Spectrum $\CSpec$}

\textbf{Correct proof: $H^1=0$.}\quad
Take the cover
$\{U_1:\text{``$p$ is even''},\;U_2:\text{``$q$ is even''}\}$.
The local proofs (even-square lemma applied to $p^2=2q^2$ and
$q^2=2k^2$) glue at the overlap
$U_1\cap U_2=\{\gcd(p,q)=1\}$, since having both $2\mid p$ and
$2\mid q$ contradicts $\gcd(p,q)=1$.  The \v{C}ech coboundary
vanishes: $\delta^0=0$, so $H^1=0$.

\textbf{Proof with a gap: $H^1\neq 0$.}\quad
Suppose we omit the hypothesis $\gcd(p,q)=1$.  Then the local
proofs ``$2\mid p$'' and ``$2\mid q$'' are both valid, but the
contradiction step has no fuel---$\gcd(p,q)$ is unconstrained.

\begin{proposition}[Gap = non-trivial $H^1$]
The proofs over $U_1$ and $U_2$ define local sections of the
proof presheaf, but $\delta^0(s_1,s_2)\neq 0$ at the overlap
$U_1\cap U_2$: the two local proofs do not compose into a
global contradiction because the necessary compatibility datum
($\gcd=1$) is missing.
\end{proposition}

Repair: adding ``without loss of generality, assume
$\gcd(p,q)=1$'' restores the missing section, making
$\delta^0=0$ and $H^1=0$.

\subsection{The Graded Evaluation $\GEval$}

\begin{align*}
  g_{\mathrm{rigor}}       &= \gone   &&\text{(every step is valid)},\\
  g_{\mathrm{elegance}}    &= -0.3  &&\text{(standard; not the most elegant proof)},\\
  g_{\mathrm{novelty}}     &= -3.0  &&\text{(2500 years old)},\\
  g_{\mathrm{significance}}&= -0.1  &&\text{(foundational result)}.
\end{align*}
Overall $G = \gone\otG(-0.3)\otG(-3.0)\otG(-0.1)=-3.0$
(bottleneck: novelty).  This is correct but unsurprising---exactly
the evaluation a journal referee would give.

\medskip

%% ─────────────────────────────────────────────────────────────
\section{Example 2: Mathematical Ideation}
\label{sec:ideation}

Consider a mathematician who notices:
$1+3=4=2^2$,\;
$1+3+5=9=3^2$,\;
$1+3+5+7=16=4^2$.

\textbf{$\STop$.}\quad
The site $\mathcal{C}_{\mathrm{conj}}$ has objects = observed
equalities and a conjectured theorem node
$T:``\sum_{i=1}^n(2i{-}1)=n^2$''.  The cover of $T$ includes
the examples plus the inductive step (which is not yet proved).

\textbf{$\PStr$.}\quad
The type is $\Pi(n:\mathbb{N}).\,\GId_{\gone}(\mathbb{N},
  \sum_{i=1}^n(2i{-}1),\,n^2)$.
A full term requires induction:
\[
  \dfrac{
    \Gamma\vdash_{\gone}\mathrm{base}:\sum_{i=1}^1(2i{-}1)=1^2
    \qquad
    \Gamma\vdash_{\gone}\mathrm{step}:
      \bigl(\textstyle\sum_{i=1}^n(2i{-}1)=n^2\bigr)
      \Rightarrow
      \bigl(\textstyle\sum_{i=1}^{n+1}(2i{-}1)=(n{+}1)^2\bigr)
  }{
    \Gamma\vdash_{\gone}\mathrm{ind}(\mathrm{base},\mathrm{step}):
      \Pi(n).\,\textstyle\sum_{i=1}^n(2i{-}1)=n^2
  }
\]
Before the inductive step is found, the mathematician has only
the examples.  These are \emph{partial evidence} at grades
$g<\gone$: three witnesses do not constitute a proof.

\textbf{$\CSpec$.}\quad
Before the inductive proof, $H^1\neq 0$: the examples cover
$n=2,3,4$ but do not glue into a proof for all~$n$.  After
finding the inductive step, $H^1=0$: the proof glues globally.
This is a precise formalisation of the difference between
\emph{conjecture} and \emph{theorem}.

\textbf{$\GEval$.}\quad
Before proof: $g_{\mathrm{rigor}}=-2.0$ (empirical only).
After proof: $g_{\mathrm{rigor}}=\gone$,
$g_{\mathrm{elegance}}=-0.1$ (clean one-line induction),
$g_{\mathrm{novelty}}=-2.5$ (well-known),
$g_{\mathrm{significance}}=-0.5$ (useful but not deep).

\begin{remark}[Ideation as proof search]
Mathematical ideation---the process of noticing a pattern and
conjecturing a theorem---is precisely the $\FMLM$ applied to
$\PStr$: generate candidate proof terms (``try induction,''
``try contradiction,'' ``try a clever substitution''), score
each against the type, and iterate.  The grade $g$ rises from
$\gzero$ (no evidence) through intermediate values (empirical
support) to $\gone$ (complete proof).  $\CSpec$ tracks which
parts of the proof are missing at each stage.
\end{remark}

\chapter{Cross-Domain Isomorphisms}
\label{ch:isomorphisms}

The preceding four chapters instantiated the 4-tuple
$(\STop,\PStr,\CSpec,\GEval)$ in four domains: programming,
natural language, poetry, and mathematics.  This chapter
demonstrates that the instantiations are not merely analogous
but \emph{functorially related}: structure-preserving maps
connect every pair.

%% ─────────────────────────────────────────────────────────────
\section{The Category of Domains}
\label{sec:dom-cat}

\begin{definition}[Domain 4-tuple]
A \emph{domain instantiation} is a quadruple
$\mathcal{D}=(\STop_{\mathcal{D}},\PStr_{\mathcal{D}},
 \CSpec_{\mathcal{D}},\GEval_{\mathcal{D}})$
where each component is the restriction of the universal
4-tuple to the domain's site, types, cohomology, and grade
components.
\end{definition}

\begin{definition}[Domain morphism]
A \emph{domain morphism} $\Phi:\mathcal{D}_1\to\mathcal{D}_2$
is a quadruple of functors/natural transformations:
\[
  \Phi = (\Phi_{\STop},\;\Phi_{\PStr},\;\Phi_{\CSpec},\;\Phi_{\GEval})
\]
such that (i)~$\Phi_{\STop}$ is a morphism of sites,
(ii)~$\Phi_{\PStr}$ preserves typing derivations,
(iii)~$\Phi_{\CSpec}$ maps cocycles to cocycles and
coboundaries to coboundaries (hence induces a map on $H^1$),
and (iv)~$\Phi_{\GEval}$ is a semiring homomorphism on grades.
\end{definition}

We write $\mathbf{Dom}$ for the category of domain instantiations
and domain morphisms.  The four domains studied form a cycle:
\[
\begin{tikzcd}
  \mathcal{D}_{\mathrm{prog}}
    \arrow[r,"\Phi_{\mathrm{prog}\to\mathrm{NL}}"]
  & \mathcal{D}_{\mathrm{NL}}
    \arrow[d,"\Phi_{\mathrm{NL}\to\mathrm{poet}}"]
  \\
  \mathcal{D}_{\mathrm{math}}
    \arrow[u,"\Phi_{\mathrm{math}\to\mathrm{prog}}"]
  & \mathcal{D}_{\mathrm{poet}}
    \arrow[l,"\Phi_{\mathrm{poet}\to\mathrm{math}}"]
\end{tikzcd}
\]

%% ─────────────────────────────────────────────────────────────
\section{Programs $\leftrightarrow$ Sentences}
\label{sec:prog-nl}

\begin{construction}[The functor $\Phi_{\mathrm{prog}\to\mathrm{NL}}$]
\label{con:prog-nl}
\begin{itemize}[nosep]
  \item $\Phi_{\STop}$: program points $\mapsto$ discourse segments;
    basic blocks $\mapsto$ phrases; the CFG cover $\mapsto$ constituent
    cover.
  \item $\Phi_{\PStr}$: specifications $\mapsto$ propositions
    (a spec ``output is sorted'' becomes the proposition ``the
    array is sorted''); programs $\mapsto$ utterances (the code
    \emph{is} a statement about what it does).
  \item $\Phi_{\CSpec}$: bugs $\mapsto$ ungrammaticality.
    A data race ($H^1\neq 0$ in the thread cover) maps to a
    binding failure ($H^1\neq 0$ in the clause cover).
  \item $\Phi_{\GEval}$: test quality $\mapsto$ linguistic quality;
    type safety $\mapsto$ grammaticality;
    performance $\mapsto$ pragmatic efficiency.
\end{itemize}
\end{construction}

\begin{theorem}[Faithfulness]
\label{thm:prog-nl-faithful}
$\Phi_{\mathrm{prog}\to\mathrm{NL}}$ is a faithful functor:
distinct programs map to distinct ``sentences'' (no two
structurally different programs have identical NL images).
\end{theorem}

\begin{proof}[Proof sketch]
Faithfulness follows from the injectivity of $\Phi_{\PStr}$ on
derivation trees: different proof trees (Hoare-logic derivations)
yield different typing derivations in $\PStr_{\mathrm{NL}}$
because each atomic Hoare rule maps to a distinct NL typing
rule.
\end{proof}

\begin{remark}[Not full]
The functor is not full: natural-language sentences expressing
emotions, social conventions, or fictional narratives have no
corresponding program.  The image of $\Phi_{\mathrm{prog}\to\mathrm{NL}}$
is the sublanguage of \emph{declarative descriptions of
computational processes}.
\end{remark}

%% ─────────────────────────────────────────────────────────────
\section{Sentences $\leftrightarrow$ Poems}
\label{sec:nl-poet}

\begin{construction}[The functor $\Phi_{\mathrm{NL}\to\mathrm{poet}}$]
\label{con:nl-poet}
\begin{itemize}[nosep]
  \item $\Phi_{\STop}$: $\mathcal{C}_{\mathrm{NL}}\hookrightarrow
    \mathcal{C}_{\mathrm{poet}}$ by adding metrical, rhyme, and
    stanzaic objects.  Every NL site embeds into the poetic site as
    the ``prose sub-site.''
  \item $\Phi_{\PStr}$: a sentence type $A$ maps to
    $A\wedge\mathrm{Meter}(\mu)\wedge\mathrm{Rhyme}(\rho)\wedge
    \mathrm{Form}(f)$---the sentence plus additional constraints.
    Every grammatical sentence is a trivial ``poem'' (free verse
    with no metre or rhyme constraints).
  \item $\Phi_{\CSpec}$: NL obstructions embed as poetic
    obstructions; poetry has \emph{additional} obstructions
    (broken metre, failed rhyme, weak volta).
  \item $\Phi_{\GEval}$: the NL grade components
    $g_{\phon},g_{\synt},g_{\sem},g_{\prag}$ embed into the
    larger poetic grade vector, which adds
    $g_{\mathrm{meter}},g_{\mathrm{rhyme}},g_{\mathrm{imagery}},
    g_{\mathrm{arc}}$.
\end{itemize}
\end{construction}

\begin{theorem}[Faithful embedding]
\label{thm:nl-poet-embed}
$\Phi_{\mathrm{NL}\to\mathrm{poet}}$ is faithful: every
grammatical sentence embeds as a (trivially constrained) poem.
The embedding is not full: poems carry structure that arbitrary
prose does not.
\end{theorem}

The key insight: poetry is \emph{NL plus extra strata in $\PStr$
and extra components in $\GEval$}.  The CTT stratal tower
$U_{\phon}:U_{\morph}:\cdots:U_{\poet}:U_{\aesth}$ makes this
precise---poetry activates the upper strata that prose leaves
dormant.

%% ─────────────────────────────────────────────────────────────
\section{Proofs $\leftrightarrow$ Programs (Curry--Howard)}
\label{sec:math-prog}

\begin{construction}[The functor $\Phi_{\mathrm{math}\to\mathrm{prog}}$]
\label{con:math-prog}
This is the classical Curry--Howard correspondence, enriched with
$\CSpec$ and $\GEval$:
\begin{itemize}[nosep]
  \item $\Phi_{\STop}$: the dependency graph of lemmas $\mapsto$
    the call graph of functions.  A cover (complete set of
    dependencies) $\mapsto$ a cover (complete set of callee
    implementations).
  \item $\Phi_{\PStr}$: propositions $\mapsto$ types;
    proofs $\mapsto$ programs; the derivation tree of a proof
    $\mapsto$ the typing derivation of a program.
  \item $\Phi_{\CSpec}$: a proof gap (missing lemma, $H^1\neq 0$)
    $\mapsto$ a bug (missing case, $H^1\neq 0$).
    Exact correspondence: non-vanishing $H^1$ in the mathematical
    site maps to non-vanishing $H^1$ in the program site.
  \item $\Phi_{\GEval}$: rigor $\mapsto$ type safety;
    elegance $\mapsto$ code cleanliness;
    novelty $\mapsto$ algorithmic innovation;
    significance $\mapsto$ practical impact.
\end{itemize}
\end{construction}

\begin{theorem}[Equivalence on $\CSpec$]
\label{thm:ch-cspec}
Under Curry--Howard,
$H^1(\mathcal{C}_{\mathrm{math}},\mathcal{F}_{\mathrm{proof}})
\cong
H^1(\mathcal{C}_{\mathrm{prog}},\mathcal{F}_{\mathrm{inv}})$.
Proof gaps and program bugs are cohomologically identical.
\end{theorem}

\begin{proof}[Proof sketch]
$\Phi_{\CSpec}$ maps the \v{C}ech complex of the lemma cover to
the \v{C}ech complex of the callee cover.  Since Curry--Howard
is an isomorphism on derivation trees, it induces an isomorphism
on cochain groups, hence on cohomology.
\end{proof}

%% ─────────────────────────────────────────────────────────────
\section{Poems $\leftrightarrow$ Proofs}
\label{sec:poet-math}

This is the most surprising correspondence.

\begin{construction}[The functor $\Phi_{\mathrm{poet}\to\mathrm{math}}$]
\label{con:poet-math}
\begin{itemize}[nosep]
  \item \textbf{Form $\mapsto$ proof structure.}\quad
    The stanzaic form of a sonnet (three quatrains + couplet)
    $\mapsto$ the logical structure of a proof (three lemmas +
    final theorem).
  \item \textbf{Meter $\mapsto$ proof rhythm.}\quad
    The alternation of stressed and unstressed syllables
    $\mapsto$ the alternation of lemma statements and their
    applications in a proof.
  \item \textbf{Rhyme $\mapsto$ structural echo.}\quad
    Rhyming end-words ($\GId_{g_r}$ at the $\phon$ stratum)
    $\mapsto$ reuse of the same lemma at different points in a
    proof (structural parallelism).
  \item \textbf{Volta $\mapsto$ key insight.}\quad
    The turn at line~9 of a sonnet $\mapsto$ the ``aha moment''
    where a proof introduces the decisive idea.
  \item \textbf{Elegance $\mapsto$ elegance.}\quad
    Both domains have an ``elegance'' component in $\GEval$ that
    resists formal definition but is recognised by experts.
\end{itemize}
\end{construction}

\begin{remark}[Why this works]
The correspondence is not metaphorical---it is structural.
Both poems and proofs are terms inhabiting a type in $\PStr$,
both must satisfy global gluing conditions ($H^1=0$ in
$\CSpec$), and both are evaluated by multi-component grades
in $\GEval$.  The functor $\Phi_{\mathrm{poet}\to\mathrm{math}}$
identifies the \emph{formal skeleton} shared by both, while
acknowledging that the \emph{content} differs (linguistic vs.\
logical).
\end{remark}

%% ─────────────────────────────────────────────────────────────
\section{The Cycle and Its Fixed Points}
\label{sec:cycle}

\begin{proposition}[The cycle is idempotent]
\label{prop:cycle-idempotent}
Let $\Psi = \Phi_{\mathrm{math}\to\mathrm{prog}}
            \circ\Phi_{\mathrm{poet}\to\mathrm{math}}
            \circ\Phi_{\mathrm{NL}\to\mathrm{poet}}
            \circ\Phi_{\mathrm{prog}\to\mathrm{NL}}$.
Then $\Psi$ is not the identity on $\mathcal{D}_{\mathrm{prog}}$
(the round trip through four domains loses domain-specific
detail), but $\Psi^2=\Psi$: the image of $\Psi$ is a
\emph{retract}, and applying $\Psi$ again does not lose further
information.
\end{proposition}

\begin{proof}[Proof sketch]
Each functor in the cycle is faithful, hence injective on
morphisms.  The composition $\Psi$ projects onto the
\emph{shared structural skeleton}---the sub-4-tuple consisting
of features present in all four domains.  This sub-4-tuple is
already in the image of $\Psi$, so $\Psi^2$ acts as the
identity on it.
\end{proof}

\begin{definition}[Universal artifact]
An artifact $a$ is \emph{universal} if $\Psi(a)\cong a$---it
is a fixed point of the cycle.  Such an artifact has coherent
realisations in all four domains simultaneously.
\end{definition}

%% ─────────────────────────────────────────────────────────────
\section{Worked Example: Recursion Across Four Domains}
\label{sec:recursion}

We take the concept of \emph{recursion} and exhibit its
4-tuple in each domain, then verify the functorial
correspondences.

\subsection{As a Program: Factorial}

\textbf{$\STop$.}\quad
The CFG contains a self-loop at the recursive call site.

\textbf{$\PStr$.}\quad
Type: $\mathrm{Factorial}:\mathbb{N}\to\mathbb{N}$,
with specification $\mathrm{fact}(n)=n!$.  The term
is $\lambda n.\;\mathrm{if}\;n=0\;\mathrm{then}\;1\;
\mathrm{else}\;n\cdot\mathrm{fact}(n-1)$.
Termination proof: the argument $n$ strictly decreases.

\textbf{$\CSpec$.}\quad
$H^1=0$: the base case and inductive case glue at $n=0$.

\textbf{$\GEval$.}\quad
$g_{\mathrm{rigor}}=\gone$, $g_{\mathrm{perf}}=-0.5$
(stack overflow for large $n$).

\subsection{As a Sentence: Centre-Embedding}

\begin{quote}
\emph{``The cat that the dog that the rat bit chased ran.''}
\end{quote}

\textbf{$\STop$.}\quad
Three nested relative clauses form a recursive phrase structure.

\textbf{$\PStr$.}\quad
The typing derivation nests three NP--VP pairs, each containing
the next.  This is linguistic recursion: the grammar rule
$\text{NP}\to\text{NP}\;\text{RelClause}$ applies to its own
output.

\textbf{$\CSpec$.}\quad
$H^1\neq 0$ for most human parsers: the triple centre-embedding
exceeds working memory, causing a processing breakdown.  The
obstruction is at the innermost clause boundary.

\textbf{$\GEval$.}\quad
$g_{\synt}=\gone$ (grammatical), $g_{\prag}=-2.0$
(incomprehensible to most listeners).

\textbf{Functor check.}\quad
$\Phi_{\mathrm{prog}\to\mathrm{NL}}$ maps the self-loop in
the CFG to the nested clause structure.  The recursive call
corresponds to the recursive application of the NP rule.

\subsection{As a Poem: A Villanelle}

The villanelle repeats two refrain lines (A1 and A2)
throughout, returning to them in a fixed pattern.  The
repetition \emph{is} recursion: the poem calls back to
its own earlier material.

\textbf{$\PStr$.}\quad
Type: $\mathrm{Villanelle}(\theta,\text{ABA}\ldots,\mu)$.
The refrain lines are shared sub-terms, reused at specified
positions---exactly like a recursive call reusing the
function body.

\textbf{$\CSpec$.}\quad
$H^1=0$ when the refrains integrate seamlessly into each
stanza.  $H^1\neq 0$ if a refrain clashes with its local
context (e.g.\ the repeated line contradicts new material).

\textbf{Functor check.}\quad
$\Phi_{\mathrm{NL}\to\mathrm{poet}}$ maps the recursive NP
structure to the refrain structure: both are self-referential
uses of material.

\subsection{As a Proof: Induction}

\textbf{$\PStr$.}\quad
Type: $\Pi(n:\mathbb{N}).\,P(n)$.  The proof term is
$\mathrm{ind}(\mathrm{base},\lambda n.\lambda h.\,
\mathrm{step}(n,h))$ where the step function receives
the inductive hypothesis~$h:P(n)$ and produces $P(n+1)$.
This \emph{is} recursion: the proof of $P(n+1)$ calls
the proof of $P(n)$.

\textbf{$\CSpec$.}\quad
$H^1=0$ when the base case and step glue at $n=0$.

\textbf{Functor check.}\quad
$\Phi_{\mathrm{math}\to\mathrm{prog}}$ (Curry--Howard)
maps the induction term to a recursive program.
$\Phi_{\mathrm{poet}\to\mathrm{math}}$ maps the
villanelle's refrain pattern to the reuse of the inductive
hypothesis.

\medskip

\begin{theorem}[Recursion as a universal artifact]
\label{thm:recursion-universal}
Recursion is a fixed point of the cycle $\Psi$: its 4-tuple
is invariant (up to isomorphism) under the round-trip
$\mathrm{prog}\to\mathrm{NL}\to\mathrm{poet}\to\mathrm{math}
\to\mathrm{prog}$.  The structural skeleton---a self-referential
term with a base case and an inductive/recursive step---is
preserved by every functor in the cycle.
\end{theorem}

\begin{remark}[Other universal artifacts]
Recursion is not the only fixed point.  Other candidates include
\emph{composition} (function composition / sentence conjunction /
stanza sequencing / lemma chaining),
\emph{abstraction} ($\lambda$-abstraction / pronoun binding /
metonymy / universal quantification), and
\emph{symmetry} (invariance under permutation / palindrome /
chiasmus / group-theoretic symmetry).  Each deserves a full
treatment parallel to the recursion example above.
\end{remark}
